% ===================================================================
%  BẢNG Số 4 — Tuổi trưởng thành và tuổi già.
%
%  Traduction, faite en 2026, en vietnamien standard d'aujourd'hui,
%  sur l'ido de texto/io. Regles de la colonne : voir 10-bang-01.tex.
%  Deux scenes, deux numerotations : les cles portent c1 et c2.
%
%  C'EST LE TABLEAU OU LA PARENTE VIETNAMIENNE SE MONTRE ENTIERE, et
%  aucune autre colonne de la serie n'a eu autant a decider. Le
%  francais dit « oncle », « tante », « beau-pere », « belle-mere »,
%  « beau-frere », « belle-soeur » sans dire de quel cote ni de quel
%  age ; le vietnamien ne PEUT pas les dire sans le dire :
%
%    ông nội   le grand-pere paternel      bà nội   la grand-mere
%    bác       le frere AINE du pere       chú      le frere CADET
%    cô        la soeur du pere            dì       la soeur de la mere
%    bố chồng  le beau-pere de l'epouse    bố vợ    celui de l'epoux
%    con dâu   la bru                      con rể   le gendre
%
%  Le repas de noces reunit les deux familles autour d'une table : il
%  a donc fallu, pour chaque personne, savoir de quel cote elle vient.
%  La gravure et le texte le disent — c'est la mere DU MARIE qui
%  sourit (17), le beau-pere DE LA MARIEE qui est heureux d'avoir une
%  bru (20) — et la colonne les nomme en consequence. La ou le livret
%  de 1926 laissait un flou, le vietnamien oblige a lire la planche.
%
%  ET LE GRAND-PERE DE LA SECONDE SCENE EST « ông nội », le paternel,
%  parce que les enfants qui viennent lui souhaiter sa fete sont ceux
%  de Gioan, son fils. Le maréchal du dernier alinea, dont le portrait
%  pend au mur, est « cụ tổ » — l'aieul, celui dont on ne compte plus
%  les degres.
%
%  LES MALADIES DE LA GRAND-MERE SONT VIETNAMIENNES, LE REMEDE EST
%  FRANCAIS : chóng mặt le vertige, cảm le rhume, đau răng le mal de
%  dents, thấp khớp le rhumatisme — mais « thuốc nước », la potion, et
%  « đơn thuốc », l'ordonnance, sont du vocabulaire de la pharmacie
%  coloniale.
% ===================================================================

%%K t04-apar-1 apar
\VUsaut{4.0mm}
\VUtitre{15.0pt}{60}{BẢNG Số 4}
\VUsaut{1.6mm}
\VUtitre{11.4pt}{0}{Tuổi trưởng thành và tuổi già.}
\VUsaut{3.2mm}

%%K t04-tit-1 sub
\VUcentre{11.4pt}{0}{\textit{Cảnh thứ nhất.}}
\VUsaut{1.6mm}
\VUtitre{11.4pt}{0}{Tiệc cưới.}
\VUsaut{2.6mm}

%%K t04-tit-2 sub
\VUpk{10.2pt}{40}{Mùa thu.}
\VUsaut{3.0mm}

%%K t04-c1-01-1 p
1. --- Các bạn trẻ đã lớn. Một trong số họ, Gioan, lấy vợ. Chúng ta dự \VUgras{tiệc
cưới}, nơi hai gia đình của đôi vợ chồng ngồi quanh cùng một bàn.

%%K t04-c1-02-1 p
2. --- Chúng ta đang ở \VUgras{mùa thu}, vào \VUgras{tháng chín}. Gió lạnh của
\VUgras{tháng mười} và \VUgras{tháng mười một} chưa lột hết màu lá xanh của đồng quê.
Không khí buổi chiều ấm áp và thơm ngát; vì thế bữa tối diễn ra dưới \VUgras{hiên nhà}
\textsuperscript{(8)}, phía vườn.

%%K t04-c1-03-1 p
3. --- Xa xa, trên \VUgras{ngọn đồi} \textsuperscript{(3)}, là nhà của cha mẹ Gioan, một
\VUgras{tòa lâu đài} \textsuperscript{(1)} thanh nhã ẩn trong màu xanh; những vườn nho
đẹp, cho thứ rượu tuyệt hảo, phủ kín sườn đồi. Người trồng nho chăm sóc chúng.

%%K t04-c1-04-1 p
4. --- Trên những vườn nho, một đàn \VUgras{gà gô} \textsuperscript{(2)} bay qua, hoảng
sợ vì \VUgras{người gác săn} \textsuperscript{(5)} đến gần. Ông này, \VUgras{súng} kẹp
dưới nách và \VUgras{túi săn} đeo chéo vai, lùng sục các \VUgras{bụi cây}
\textsuperscript{(4)} cùng con \VUgras{chó} \textsuperscript{(6)} của mình. Ông trông coi
điền trang và ngăn không cho kẻ săn trộm hủy hoại thú rừng.

%%K t04-c1-05-1 p
5. --- Bên ngoài, dựa vào \VUgras{hiên nhà}, một giàn nho leo lên, buông xuống những
\VUgras{chùm nho} \textsuperscript{(7)} nặng trĩu.

%%K t04-c1-06-1 p
6. --- Những bông hoa mùa thu đẹp đẽ trang trí \VUgras{bàn tiệc} \textsuperscript{(16)}
là \VUgras{hoa cúc} \textsuperscript{(9)}; \VUgras{cha} \textsuperscript{(10)} của cô dâu
đã cài một bông lên khuyết áo lễ phục của mình.

%%K t04-c1-07-1 p
7. --- Bữa ăn sắp tàn. \VUgras{Người hầu bàn} \textsuperscript{(11)} bắt đầu mang các
món tráng miệng ra và lát nữa sẽ dọn đi những thứ trong bộ đồ ăn không còn cần đến :
\VUgras{thìa} \textsuperscript{(14)}, \VUgras{nĩa} \textsuperscript{(12)},
\VUgras{dao} \textsuperscript{(13)} và \VUgras{đĩa lớn} \textsuperscript{(15)}.

%%K t04-tit-3 sub
\VUpk{10.2pt}{40}{Họ hàng.}
\VUsaut{3.0mm}

%%K t04-c1-08-1 p
8. --- Ai nấy đều có vẻ vui; và nụ cười của \VUgras{mẹ} \textsuperscript{(17)} chú rể
khi nhìn các con mình cho thấy bà hạnh phúc.

%%K t04-c1-09-1 p
9. --- Đám cưới này kết hợp hai gia đình giàu có trong vùng. Gioan, \VUgras{chú rể}
\textsuperscript{(18)}, là \VUgras{con trai} của một điền chủ lớn và trở thành
\VUgras{con rể} của một nhà công nghiệp tài giỏi.

%%K t04-c1-10-1 p
10. --- Đôi vợ chồng còn trẻ, thế mà họ đã \VUgras{đính hôn} từ lâu. \VUgras{Cô dâu}
\textsuperscript{(19)}, Marta, thật duyên dáng trong \VUgras{áo cưới trắng} cài hoa cam.

%%K t04-c1-11-1 p
11. --- \VUgras{Bố chồng} \textsuperscript{(20)} của cô rất sung sướng có được một
\VUgras{con dâu} hiền hậu như thế, còn \VUgras{mẹ vợ} \textsuperscript{(21)} của Gioan
mỉm cười khi thấy \VUgras{con gái} mình đẹp đến vậy.

%%K t04-c1-12-1 p
12. --- Ở cuối bàn ngồi những \VUgras{khách mời} \textsuperscript{(22)} khác :
\VUgras{họ hàng, bạn bè, anh em họ, chị em họ}. Một \VUgras{người quản tiệc}
\textsuperscript{(23)}, khăn ăn vắt trên tay, ân cần phục vụ họ.

%%K t04-tit-4 sub
\VUpk{10.2pt}{40}{Bạn bè.}
\VUsaut{3.0mm}

%%K t04-c1-13-1 p
13. --- Bên phải bàn tiệc, kìa một \VUgras{cô gái} \textsuperscript{(24)},
\VUgras{bạn} của cô dâu, rồi \VUgras{anh} của cô ấy, người làm \VUgras{phù rể}
\textsuperscript{(25)}. Anh trò chuyện với \VUgras{cô phù dâu} \textsuperscript{(26)},
em gái chú rể, người nhờ đám cưới này mà thành \VUgras{em chồng} của Marta. Đó là một
thiếu nữ duyên dáng. Cô có những món \VUgras{nữ trang} đẹp, một \VUgras{chuỗi ngọc
trai} và một chiếc \VUgras{vòng tay} bằng vàng.

%%K t04-c1-14-1 p
14. --- Gần họ, người đàn ông đứng lên và nâng ly là một \VUgras{người bạn}
\textsuperscript{(27)} của gia đình. Ông là một trong những \VUgras{người làm chứng cho
chú rể}. Ông \VUgras{nâng cốc chúc mừng}, uống mừng \VUgras{sức khỏe} đôi tân hôn và
chúc họ hạnh phúc dài lâu. Ông mặc lễ phục : áo đuôi tôm và \VUgras{giày da bóng}
\textsuperscript{(28)}. Chính ông đưa \VUgras{bà nội} \textsuperscript{(29)}, người đã
\VUgras{góa chồng}.

%%K t04-c1-15-1 p
15. --- Chàng trai mặc \VUgras{áo đuôi tôm} \textsuperscript{(31)}, có bộ ria mép đen
mảnh, là \VUgras{anh rể} \textsuperscript{(30)} của Gioan. Anh là một kỹ sư xuất sắc.
Anh ngồi cạnh một \VUgras{thiếu phụ} \textsuperscript{(33)} rất thanh lịch, \VUgras{chị
cả} của Marta và là mẹ của cô bé xinh xắn đang khoác tay anh họ đến tặng hoa và
\VUgras{chúc} bà \VUgras{buổi tối tốt lành}. Bà ấy đeo đôi \VUgras{hoa tai} rất đẹp và
có \VUgras{kiểu tóc} thanh lịch.

%%K t04-c1-16-1 p
16. --- Cậu bé anh họ, trông thật ngộ với chiếc \VUgras{áo choàng}
\textsuperscript{(36)} và chiếc \VUgras{quần chẽn} \textsuperscript{(37)} phồng, thật
đáng thương. Em là \VUgras{trẻ mồ côi} \textsuperscript{(35)}. Em mất cha mẹ từ rất nhỏ.
Bác em là \VUgras{người giám hộ} \textsuperscript{(38)} của em và đã ở vậy không lấy vợ
để nuôi đứa bé đáng thương. Ông là người tốt bụng. Ông hơi hói và rất cận thị; ông đeo
\VUgras{kính kẹp mũi}.

%%K t04-tit-5 sub
\VUsaut{2.4mm}
\VUpk{10.2pt}{40}{Món tráng miệng.}
\VUsaut{3.0mm}

%%K t04-c1-17-1 p
17. --- Sau lưng thực khách, trên một chiếc bàn phủ \VUgras{khăn}, \VUgras{món tráng
miệng} \textsuperscript{(39)} đã được dọn sẵn. Trong một chiếc cốc lớn là trái cây :
\VUgras{đào} \textsuperscript{(40)}, \VUgras{lê} \textsuperscript{(41)}, \VUgras{táo}
\textsuperscript{(42)}, \VUgras{mơ} \textsuperscript{(43)} và \VUgras{nho}
\textsuperscript{(44)}. Bên cạnh là \VUgras{dứa} \textsuperscript{(45)}, một chiếc
\VUgras{bánh ngọt} \textsuperscript{(46)} đẹp, những \VUgras{chai rượu mạnh}
\textsuperscript{(48)} và \VUgras{rượu sâm banh} \textsuperscript{(49)}, cùng những chồng
\VUgras{đĩa} \textsuperscript{(47)} sạch.

%%K t04-c1-18-1 p
18. --- \VUgras{Người hầu rượu} \textsuperscript{(50)} mở một \VUgras{chai}
\textsuperscript{(52)} rượu vang lâu năm bằng \VUgras{cái mở nút}
\textsuperscript{(51)}. \VUgras{Nút bần} \textsuperscript{(53)} có vẻ rất khó rút. Người
hầu rượu ấy vắt một chiếc \VUgras{khăn} \textsuperscript{(54)} trên tay.

%%K t04-c1-19-1 p
19. --- Sau bữa tối, các ông sẽ sang phòng hút thuốc, còn các bà sẽ đi sửa soạn cho
buổi khiêu vũ.

%%K t04-tit-6 sub
\VUsaut{5.0mm}
\VUcentre{11.4pt}{0}{\textit{Cảnh thứ hai.}}
\VUsaut{1.6mm}
\VUpk{11.4pt}{40}{Sinh nhật của ông nội.}
\VUsaut{2.6mm}

%%K t04-tit-7 sub
\VUpk{10.2pt}{40}{Cuộc thăm viếng.}
\VUsaut{3.0mm}

%%K t04-c2-01-1 p
1. --- Mười năm đã qua. Cha mẹ Gioan đã già và rất yêu ba đứa cháu của mình, hai
\VUgras{cháu trai} \textsuperscript{(2)} và một \VUgras{cháu gái}
\textsuperscript{(3)}. Vì hôm nay là \VUgras{sinh nhật của ông nội}, các con đến chúc
mừng ông.

%%K t04-c2-02-1 p
2. --- Bé Phêrô còn rất nhỏ, em mới ba tuổi. Em thấy con \VUgras{mèo}
\textsuperscript{(1)} lại gần. Em muốn vuốt bộ \VUgras{lông} mượt như tơ và cái
\VUgras{đuôi} dài của nó; em không nghĩ rằng dưới những cái \VUgras{chân} duyên dáng ấy
có giấu những chiếc vuốt nhọn hiểm ác có thể cào em.

%%K t04-c2-03-1 p
3. --- Chị em, Gabriela, người đang nắm tay em, nghiêm nghị hơn nhiều; còn Marcelô, với
chiếc \VUgras{áo khoác} \textsuperscript{(4)} rộng và chiếc \VUgras{thắt lưng}
\textsuperscript{(5)} da, trông ra dáng người lớn, dù chiếc quần ngắn để lộ đôi
\VUgras{tất} \textsuperscript{(6)} của em.

%%K t04-c2-04-1 p
4. --- Cha các em mặc chiếc \VUgras{áo bành tô} \textsuperscript{(7)} mùa đông dày có
\VUgras{cổ lông} \textsuperscript{(8)}, và trong \VUgras{túi} \textsuperscript{(9)} có
chiếc khăn quàng. Vợ ông đội mũ dạ cài \VUgras{lông chim} \textsuperscript{(10)}, mặc
\VUgras{áo vét} \textsuperscript{(11)}, quàng \VUgras{khăn lông} \textsuperscript{(12)}
và mang \VUgras{bao tay lông} \textsuperscript{(13)}.

%%K t04-c2-05-1 p
5. --- Trời rất lạnh, vì đang là mùa đông. \VUgras{Tuyết} \textsuperscript{(19)} đã rơi
suốt ngày và mái nhà trắng xóa. Cùng với \VUgras{đêm}, cái lạnh càng thêm buốt. Trên
trời, \VUgras{mặt trăng} \textsuperscript{(17)} và những \VUgras{ngôi sao}
\textsuperscript{(18)} chiếu sáng, xuyên qua \VUgras{bóng tối}.

%%K t04-tit-8 sub
\VUsaut{4.2mm}
\VUpk{10.2pt}{40}{Bệnh tật.}
\VUsaut{3.0mm}

%%K t04-c2-06-1 p
6. --- \VUgras{Người mù} \textsuperscript{(20)} đáng thương đang băng qua quảng trường
trước \VUgras{hiệu thuốc} \textsuperscript{(16)}, do con \VUgras{chó xù}
\textsuperscript{(21)} của mình dẫn đường, thật là bất hạnh. Chị Iustina, \VUgras{người
giúp việc} \textsuperscript{(14)}, mang từ bếp lên một \VUgras{bát}
\textsuperscript{(15)} \VUgras{nước thuốc} thật nóng, pha nhiều đường.

%%K t04-c2-07-1 p
7. --- \VUgras{Bà nội} \textsuperscript{(23)}, muốn đi lấy \VUgras{kính}
\textsuperscript{(26)} của mình trên bàn để đọc \VUgras{đơn thuốc}
\textsuperscript{(27)} mà \VUgras{bác sĩ} \textsuperscript{(25)} vừa viết, đứng dậy khỏi
ghế bành, chống lên chiếc \VUgras{gậy} \textsuperscript{(24)} của mình.

%%K t04-c2-08-1 p
8. --- Bà thường bị những cơn \VUgras{khó ở}, \VUgras{chóng mặt}, \VUgras{cảm} và
\VUgras{đau răng}; chân bà yếu và mắt bà không còn tinh nữa. Dù vậy, bà vẫn thích trêu
ông \VUgras{bác sĩ} già của mình, người lúc nào cũng muốn kê cho bà một thang
\VUgras{thuốc nước} \textsuperscript{(28)}. Hôm nay bà rất vui vì có cả gia đình trẻ
quanh mình.

%%K t04-c2-09-1 p
9. --- Trong căn phòng đóng kín thật dễ chịu. Trong \VUgras{lò sưởi}
\textsuperscript{(29)} bằng đá hoa, một \VUgras{ngọn lửa đẹp} \textsuperscript{(30)}
bùng cháy, và người ta thấy hạnh phúc dưới \VUgras{ánh sáng} dịu của chiếc \VUgras{đèn}
\textsuperscript{(31)} vừa được thắp lên.

%%K t04-tit-9 sub
\VUsaut{4.2mm}
\VUpk{10.2pt}{40}{Ông nội.}
\VUsaut{3.0mm}

%%K t04-c2-10-1 p
10. --- Chính \VUgras{ông nội} \textsuperscript{(33)} cũng quên rằng mình đang ốm, rằng
bệnh \VUgras{thấp khớp} giữ chặt ông trong chiếc \VUgras{ghế bành}
\textsuperscript{(34)}, và rằng mới hôm qua ông còn \VUgras{sốt} và \VUgras{ngất}. Ông
không còn nhớ rằng mùa đông năm ngoái ông bị \VUgras{viêm phổi} và một cơn \VUgras{ho}
khan, nhọc nhằn đã hành hạ ông.

%%K t04-c2-10-2 p
Vậy mà ông đã rất khó khăn mới khỏi bệnh. Từ đó ông đỡ hơn một chút. Nhưng chân ông, đặt
trên chiếc \VUgras{gối} \textsuperscript{(38)} kê trên chiếc \VUgras{ghế đẩu}
\textsuperscript{(39)} nhỏ, cũng vẫn làm ông đau.

%%K t04-c2-11-1 p
11. --- Khi ông được quấn cẩn thận trong chiếc \VUgras{áo choàng trong nhà}
\textsuperscript{(35)} ấm áp có những chiếc \VUgras{cúc} \textsuperscript{(36)} lớn bằng
xà cừ, và khi bên cạnh ông có cô bé \VUgras{mồ côi} \textsuperscript{(40)} đội
\VUgras{mũ vải} trắng, mặc \VUgras{áo} \textsuperscript{(42)} xanh và \VUgras{áo choàng
ngắn} \textsuperscript{(41)} đọc báo cho ông nghe, thì ông không than phiền gì nữa.

%%K t04-c2-12-1 p
12. --- Mỗi năm, khi \VUgras{tờ lịch} \textsuperscript{(43)} lại chỉ \VUgras{ngày} mồng
một \VUgras{tháng chạp}, ông sốt ruột chờ các \VUgras{cháu} của mình, vì ông biết chúng
sẽ không quên \VUgras{ngày} quan trọng ấy.

%%K t04-c2-13-1 p
13. --- Ông xúc động nhớ lại thời rất xa xưa, khi chính ông đến chúc mừng vị
\VUgras{nguyên soái} lừng lẫy của Đế chế, người có \VUgras{bức chân dung}
\textsuperscript{(44)} treo trên tường và là \VUgras{cụ tổ} của các con ông.
\VUsaut{6.0mm}
\VUfilet{22mm}
